%=============================================================================
% W4i21: POST-W'(0)=0 NEW PREDICTIONS + PREDICTION AUDIT UPDATE
% Agent 15 (Predictions + beta=0 Follow-up) | Opus 4.6 | 20 March 2026
%
% PARADIGM STATE (i20, AUTHOR-CONFIRMED):
%   - W'(0) = 0 confirmed. Simple mu(x) = x/(1+x) requires it.
%   - Deep-MOND G_eff = G_N * sqrt(a_0/g_N) >> G_N is CORRECT.
%   - Growth delta ~ a^2 in deep-MOND (twice CDM rate).
%   - Linearized perturbation theory DEGENERATE at mu(0)=0.
%   - Nonlinear AQUAL perturbation theory is the correct framework.
%   - beta = 0 at ~2-3sigma (not 7sigma). SO decisive 2027-2028.
%   - Prediction audit: 9A, 10B, 6C. Honest 19:1.
%   - Cosmological sector RESCUED by deep-MOND quasi-static growth.
%   - sigma_8 ~ 0.5-2 (viable, may overshoot).
%   - mochi_class with nonlinear AQUAL solver: $15-20K, 3-4 weeks.
%
% MANDATE:
%   1. New predictions from nonlinear cosmological perturbation theory.
%   2. Updated prediction audit (grades, sigma_8 sharpness).
%   3. Stiff-field signatures survival check under W'(0)=0.
%   4. Sharpest near-term prediction for Euclid DR1 / DESI DR2.
%   5. DFD vs superfluid DM vs AeST comparison with deep-MOND cosmology.
%=============================================================================

\documentclass[11pt]{article}
\usepackage{amsmath,amssymb,amsthm,booktabs,geometry,xcolor,tcolorbox,longtable}
\geometry{margin=1in}

\newcommand{\DFD}{DFD}
\newcommand{\LCDM}{$\Lambda$CDM}
\newcommand{\CS}{Chern--Simons}
\newcommand{\Geff}{G_{\rm eff}}
\newcommand{\aM}{a_0}
\newcommand{\astar}{a_*}

\newtheorem{theorem}{Theorem}
\newtheorem{lemma}[theorem]{Lemma}
\newtheorem{proposition}[theorem]{Proposition}
\newtheorem{corollary}[theorem]{Corollary}
\newtheorem{finding}{Finding}
\newtheorem{prediction}{Prediction}

\title{\textbf{W4i21: Post-$W'(0)=0$ New Predictions,\\
Deep-MOND Cosmological Signatures, and\\
the Sharpest Near-Term DFD Tests}\\[6pt]
{\large Nonlinear Perturbation Theory, Prediction Audit Update,\\
and Theory Comparison}}
\author{DFD Research Programme --- W4i21 Agent 15 (Predictions)\\
Opus 4.6}
\date{20 March 2026}

\begin{document}
\maketitle

%=============================================================================
\section*{Executive Summary}
%=============================================================================

\begin{tcolorbox}[colback=blue!5!white, colframe=blue!70!black,
  title=\textbf{TOP 5 FINDINGS (Ranked by Impact on DFD Programme)}]

\begin{enumerate}

\item \textbf{FINDING 1 (Scale-Dependent Growth Rate: The Sharpest
  Near-Term DFD Prediction).}
  With $W'(0) = 0$ confirmed and the deep-MOND nonlinear AQUAL
  perturbation theory established, DFD predicts a \emph{qualitatively
  different} growth rate $f\sigma_8(k,z)$ from \LCDM.  In \LCDM, the
  growth rate $f(z) = d\ln\delta/d\ln a$ is scale-independent at
  linear order: $f \approx \Omega_m(z)^{0.55}$.  In DFD, the
  deep-MOND $\Geff = G_N\sqrt{\aM/g_N}$ depends on the Newtonian
  acceleration $g_N \propto k\,\delta_b$, making $f$ both
  scale-dependent and amplitude-dependent:
  \begin{equation}
  f^{\rm DFD}(k,z) \approx 2 - \frac{1}{4}\frac{d\ln g_N}{d\ln a}
  \quad\text{(deep-MOND limit, } g_N \ll \aM\text{)}.
  \end{equation}
  In the pure deep-MOND regime where $\delta \propto a^2$,
  $f = d\ln\delta/d\ln a = 2$.  This is dramatically larger than the
  \LCDM{} value $f \approx 0.5$ at $z = 0$ and $f \approx 1$ at high
  redshift.  However, the MOND-to-Newtonian crossover is
  scale-dependent: large $k$ (small scales, higher $g_N$) cross into
  the Newtonian regime ($\Geff \to G_N$, $f \to 0.3$) earlier than
  small $k$ (large scales).  \textbf{DFD therefore predicts a
  characteristic $f\sigma_8$ that increases toward large scales
  (small $k$)}, the opposite of any standard dark energy or modified
  gravity model operating within the Horndeski framework.

  \textbf{Testable:} Euclid DR1 (October 2026) will measure
  $f\sigma_8$ in $\sim$4 redshift bins with $\sim$5--10\% precision.
  The DFD prediction of $f\sigma_8$ \emph{increasing} toward larger
  scales at fixed redshift is a zero-parameter, qualitative
  discriminant.  Any $>3\sigma$ detection of scale-dependent growth
  with the predicted sign would be strong evidence for MOND-type
  gravity.  A null detection (scale-independent $f\sigma_8$ consistent
  with \LCDM) would constrain the deep-MOND regime to apply only at
  scales $k < k_{\rm data}$.

  \textbf{Derivation chain:} $W'(0)=0$ (i20 Finding~1, Eq.~27)
  $\Rightarrow$ linearized PT degenerate (Theorem~1)
  $\Rightarrow$ nonlinear AQUAL (Nusser 2002)
  $\Rightarrow$ $\Geff(k,\delta_b)$ from i20 Eqs.~(22)--(24)
  $\Rightarrow$ baryon continuity + Euler equations
  $\Rightarrow$ scale-dependent $f(k,z)$.

\item \textbf{FINDING 2 (Modified Power Spectrum Shape: Blue Tilt at
  Large Scales from $\Geff \propto 1/\sqrt{k\delta_b}$).}
  In \LCDM, the matter power spectrum $P(k)$ has a characteristic
  ``turnover'' at $k_{\rm eq} \approx 0.01\,h\,{\rm Mpc}^{-1}$
  (matter-radiation equality), with $P(k) \propto k$ for $k \ll k_{\rm eq}$
  and $P(k) \propto k^{-3}\ln^2(k)$ for $k \gg k_{\rm eq}$ (Harrison--Zel'dovich
  $\to$ Meszaros).  The CDM potential wells are established during radiation
  domination and baryons fall in after decoupling.

  In DFD, there is no CDM.  Structure formation is driven entirely
  by baryonic perturbations under $\Geff$.  The deep-MOND enhancement
  $\Geff \propto 1/\sqrt{k\delta_b}$ means:
  \begin{itemize}
  \item \textbf{Large-scale modes ($k \lesssim 0.01\,h\,{\rm Mpc}^{-1}$)
    are MORE enhanced} (smaller $g_N$, deeper into MOND regime).
  \item \textbf{Small-scale modes ($k \gtrsim 0.1\,h\,{\rm Mpc}^{-1}$)
    are LESS enhanced} (larger $g_N$, closer to Newtonian).
  \end{itemize}
  The net effect is a \textbf{blue tilt} of $P(k)$ at large scales
  relative to \LCDM, and a relative suppression at small scales.
  The matter-radiation equality feature in $P(k)$ is \textbf{weakened
  but not absent}: it is imprinted in the baryon transfer function at
  decoupling (through acoustic oscillations), but the post-decoupling
  growth is no longer scale-independent.

  \textbf{Key observable:} The ratio $P^{\rm DFD}(k)/P^{\LCDM}(k)$ is
  a monotonically decreasing function of $k$ in the range
  $0.001 < k < 1\,h\,{\rm Mpc}^{-1}$.  This is qualitatively similar
  to warm dark matter suppression at small scales, but the mechanism
  (enhanced large-scale growth, not free-streaming cutoff) and the
  shape (power-law tilt, not exponential cutoff) differ.

  \textbf{Quantification requires mochi\_class.}  The amplitude of the
  tilt depends on the crossover scale $k_\mu$ where $g_N(k) = \aM$,
  which depends on the nonlinear growth history.  Order-of-magnitude:
  $P^{\rm DFD}/P^{\LCDM}$ could range from $\sim$2--5 at $k = 0.01$
  to $\sim$0.3--0.8 at $k = 0.3$, but these numbers are
  \emph{conjectural} until the Boltzmann code runs.

  \textbf{Grade: C (Conjecture).}  Qualitative prediction is robust;
  quantitative amplitude requires numerical computation.

\item \textbf{FINDING 3 (Stiff-Field Signatures SURVIVE $W'(0)=0$:
  Gravitational Wave Background Unaffected).}
  The i19 stiff-field signatures derive from the \emph{temporal}
  sector of the DFD action: the free-wave equation
  $\ddot{\delta\psi}/c^2 = 0$ (no spatial gradient term at leading
  order for deep-MOND background, confirmed by i20 Eq.~(29)),
  i.e., the temporal kinetic term $(\dot{\psi})^2/c^2$ is
  \emph{always} present regardless of $W'(0)$.

  The three stiff-field signatures from i19:
  \begin{enumerate}
  \item[(a)] \textbf{Enhanced high-$k$ GW background from stiff era:}
    The $\psi$-field with $w = 1$ (stiff, $c_s = c$) in the early
    universe dilutes as $\rho_\psi \propto a^{-6}$ when it dominates.
    This amplifies primordial GWs at frequencies that re-entered the
    horizon during any stiff-dominated epoch.  The signature is a
    \emph{blue-tilted} stochastic GW background $\Omega_{\rm GW}(f)
    \propto f^2$ above the standard scale-invariant inflationary spectrum.
    \textbf{Survives.}  The $W'(0) = 0$ correction means the spatial
    kinetic term vanishes at linear order, but the stiff-era GW
    enhancement depends on the background equation of state (temporal
    sector), not perturbation gradients.  The GW tensor modes decouple
    from the scalar sector entirely.

  \item[(b)] \textbf{Absence of $\psi$-isocurvature modes:}
    CDM can carry primordial isocurvature perturbations (CDM density
    fluctuations uncorrelated with radiation).  DFD's $\psi$-field
    with $c_s = c$ has its perturbations smoothed out on all sub-Hubble
    scales, so it cannot carry isocurvature modes at recombination.
    \textbf{Survives.}  The $W'(0) = 0$ correction makes the smoothing
    even more extreme: without a spatial gradient restoring force at
    linear order, perturbations in $\psi$ itself are not propagated
    (they are instead driven by matter through the nonlinear AQUAL
    equation).  The absence of isocurvature is strengthened, not weakened.

  \item[(c)] \textbf{Jeans-scale cutoff at $k_J \sim H/c$:}
    This i19 prediction needs revision.  With $W'(0) = 0$, there is
    no linear Jeans scale for $\psi$ perturbations (the linear
    equation is degenerate).  The effective ``Jeans scale'' is set by
    the MOND crossover: scales where $g_N \sim \aM$ transition from
    deep-MOND to Newtonian behaviour.  \textbf{Modified:} the cutoff is
    not at $k_J \sim H/c$ but at the scale where the nonlinear growth
    transitions from $\delta \propto a^2$ to $\delta \propto a^{0.3}$.
    This is the MOND crossover scale $k_\mu(z)$, which evolves with
    redshift and depends on the matter density.
  \end{enumerate}

  \textbf{Net result: 2 of 3 stiff-field signatures survive intact;
  the third is modified but the physical mechanism (MOND crossover
  scale replacing Jeans scale) is well-defined.}

  \textbf{Grade: (a) B (Derivation), (b) A (Theorem), (c) C (Conjecture,
  pending mochi\_class quantification of $k_\mu(z)$).}

\item \textbf{FINDING 4 (Prediction Audit Update: 3 Predictions
  Change Grade; $\sigma_8$ Becomes LESS Sharp).}

  With $W'(0) = 0$ confirmed and deep-MOND growth established, the
  prediction audit (i19: 9A, 10B, 6C) is updated:

  \begin{table}[h]
  \centering
  \caption{Predictions with changed grades after i20/i21.}
  \begin{tabular}{llll}
  \toprule
  \textbf{Prediction} & \textbf{i19 Grade} & \textbf{i21 Grade}
    & \textbf{Reason} \\
  \midrule
  P25: $f\sigma_8(k,z)$ & (new) & B (Derivation)
    & Scale-dependent growth from AQUAL. \\
  & & & Qualitative sign derived; \\
  & & & amplitude requires mochi\_class. \\[3pt]
  P6: $\sigma_8$ value & B & C (Conjecture)
    & DOWNGRADED. Deep-MOND gives \\
  & & & $\sigma_8 \sim 0.5$--$2$, a factor \\
  & & & $\sim$4 range. No longer a sharp \\
  & & & prediction without mochi\_class. \\[3pt]
  P26: $P(k)$ blue tilt & (new) & C (Conjecture)
    & Qualitative sign robust, \\
  & & & amplitude conjectural. \\[3pt]
  P3: BAO amplitude & B & B (unchanged)
    & BAO is geometric (distance-bias \\
  & & & framework). Growth affects \\
  & & & broadband P(k), not BAO wiggles. \\
  \bottomrule
  \end{tabular}
  \end{table}

  \textbf{Updated totals:} 9A, 9B, 8C. Honest (A+B):parameter
  ratio = 18:1 (unchanged, as new P25 is B grade but P6 dropped
  from B to C).

  \textbf{The $\sigma_8$ prediction is LESS sharp.}  Before i20,
  the cosmological sector was either ``dead'' (no structure formation)
  or ``fixed'' (some specific growth rate).  After i20, the sector is
  ``viable but wide'': $\sigma_8 \sim 0.5$--$2$ depending on the
  nonlinear growth history.  The deep-MOND $\delta \propto a^2$
  growth rate is so fast that the real constraint will be
  \emph{overshooting} $\sigma_8$ (structures form too early and become
  too nonlinear), not undershooting.  The early onset of nonlinearity
  ($z \sim 5$--$10$) compared to \LCDM{} ($z \sim 1$--$3$) may actually
  \textbf{help} with the ``impossibly early galaxy'' problem
  (JWST high-$z$ galaxies), but risks conflict with the Lyman-$\alpha$
  forest power spectrum at $z \sim 2$--$4$.

  \textbf{The SHARPEST near-term prediction is NOT $\sigma_8$ but
  the scale-dependence of $f\sigma_8$} (Finding 1).  The qualitative
  sign of the scale-dependence is robust and model-independent within
  the DFD framework.

\item \textbf{FINDING 5 (DFD vs Superfluid DM vs AeST: The Key
  Distinguishing Observables with Deep-MOND Cosmology).}

  With the deep-MOND cosmological sector now established, DFD's
  landscape position relative to the two main competitors can be
  sharpened:

  \begin{table}[h]
  \centering
  \caption{Theory comparison: DFD vs Superfluid DM vs AeST (relativistic MOND).}
  \small
  \begin{tabular}{p{3.2cm}p{3.5cm}p{3.5cm}p{3.5cm}}
  \toprule
  \textbf{Observable} & \textbf{DFD} & \textbf{Superfluid DM}
    & \textbf{AeST (Skordis--Z\l{}o\'snik)} \\
  \midrule
  Dark matter particle & None (scalar field $\psi$)
    & Superfluid phonon + normal-fluid DM
    & Vector + scalar fields \\[3pt]
  GW gravitational lensing
    & \textbf{NEVER} (structural theorem, Tier 0)
    & Standard (has DM halo)
    & Standard (metric coupling) \\[3pt]
  $c_T$ (GW speed) & $= c$ exactly
    & $= c$ (GR tensor sector)
    & $= c$ (constrained) \\[3pt]
  Cosmic birefringence $\beta$
    & $= 0$ exactly (microsector)
    & No prediction
    & No prediction \\[3pt]
  CMB power spectrum
    & Standard (same $H(z)$, $\psi$ smooth)
    & Tuned to match (2 extra parameters)
    & Fits CMB (2 extra parameters) \\[3pt]
  $f\sigma_8$ scale dependence
    & \textbf{YES} (increasing toward large scales)
    & NO (standard CDM-like growth in normal fluid)
    & Possible (depends on vector field mass) \\[3pt]
  $P(k)$ shape
    & Blue-tilted (large-scale enhancement)
    & Standard CDM-like + small-scale phonon effects
    & Standard CDM-like \\[3pt]
  Vortices in galaxy haloes
    & \textbf{NO} (classical scalar, no topology)
    & \textbf{YES} (quantized vortex array, Berezhiani--Khoury 2025)
    & NO \\[3pt]
  Cluster collisions
    & Baryon-follows-gas always
    & Superfluid/normal-fluid separation
    & Normal DM behaviour \\[3pt]
  Free parameters
    & 1 ($\lambda$); $\mu(x)$ derived
    & 3--4 (DM mass, $\Lambda$, $\alpha$, $\beta$)
    & 2--3 ($\mu$-function shape, coupling constants) \\[3pt]
  $D_L^{\rm EM}/D_L^{\rm GW}$
    & $e^{\Delta\psi(z)} \neq 1$ (1.315 at $z=1$)
    & $= 1$
    & $= 1$ \\
  \bottomrule
  \end{tabular}
  \end{table}

  \textbf{The single most powerful DFD-vs-competitors discriminant is
  still T0 (GWs never gravitationally lensed)}, testable with a single
  multiply-lensed GW event from ET/CE ($\sim$2036--2040).

  \textbf{The NEAREST discriminant (2026--2027) is scale-dependent
  $f\sigma_8$.}  Both superfluid DM and AeST produce CDM-like
  scale-independent growth at the linear level (their DM components
  cluster normally).  Only DFD, with its bare-baryon deep-MOND growth,
  predicts the characteristic scale-dependent $f\sigma_8$ described in
  Finding~1.  Euclid DR1 (October 2026) can test this.

  \textbf{Caution:}  AeST's vector field \emph{can} produce
  scale-dependent growth if the vector mass is in the right range
  (Skordis \& Z\l{}o\'snik 2021, PRD).  But the mechanism differs:
  AeST's scale dependence arises from the vector field's own
  clustering, not from nonlinear MOND enhancement.  The two signatures
  have different $k$-dependence and $z$-dependence.  Disentangling
  them requires the Boltzmann codes for both theories.

  \textbf{Secondary discriminant: $D_L^{\rm EM}/D_L^{\rm GW}$} from
  bright sirens.  DFD's $\psi$-screen predicts $D_L^{\rm EM} > D_L^{\rm GW}$
  by a factor that grows with redshift ($e^{\Delta\psi(z)}$, reaching
  1.315 at $z = 1$).  Neither superfluid DM nor AeST predicts any
  EM-GW luminosity distance discrepancy.  This requires $\sim$25
  bright siren events at $z > 0.3$ (ET+CE era, $\sim$2036+).

\end{enumerate}
\end{tcolorbox}


%=============================================================================
\section{New Predictions from Nonlinear Cosmological Perturbation Theory}
\label{sec:new-predictions}
%=============================================================================

The confirmation of $W'(0) = 0$ and the deep-MOND $\Geff$ resolution
(i20) establish a qualitatively new regime for DFD cosmological
predictions.  We identify five new or sharpened predictions.

\subsection{P25: Scale-Dependent Growth Rate $f\sigma_8(k,z)$}

\begin{prediction}[Scale-Dependent Growth Rate]
\label{pred:fsigma8}
DFD predicts that the linear growth rate $f(z) \equiv d\ln\delta_b/d\ln a$
is \textbf{scale-dependent}, with $f$ increasing toward larger scales
(smaller $k$).  In the deep-MOND limit:
\begin{equation}
f^{\rm DFD}(k,z) \approx 2 \quad (k \ll k_\mu(z)),
\qquad
f^{\rm DFD}(k,z) \to f_{\rm baryon}(z) \approx 0.3 \quad (k \gg k_\mu(z)),
\end{equation}
where $k_\mu(z)$ is the MOND crossover wavenumber at which
$g_N(k,z) = \aM$.
\end{prediction}

\paragraph{Physical mechanism.}
The deep-MOND $\Geff = G_N\sqrt{\aM/g_N}$ with
$g_N = 4\pi G_N \bar\rho_b \delta_b / (k/a)$ means:
\begin{equation}
\Geff(k,z,\delta_b) = G_N \sqrt{\frac{\aM\,k}{4\pi G_N a\,\bar\rho_b(z)\,\delta_b(k,z)}}.
\end{equation}
Large-scale modes (small $k$) have smaller $g_N$ and therefore larger
$\Geff$, driving faster growth.  The growth exponent transitions from
$p \approx 2$ (deep-MOND) to $p \approx 0.3$ (baryon-only Newtonian)
as $k$ increases through $k_\mu$.

\paragraph{Crossover scale estimate.}
Setting $g_N = \aM$:
\begin{equation}
k_\mu(z) \sim \frac{\aM\,a}{4\pi G_N \bar\rho_b(z)\,\delta_b},
\end{equation}
which depends on the density contrast $\delta_b$ itself (the system is
nonlinear).  At $z = 0$ with $\bar\rho_b \sim 4 \times 10^{-28}\,
{\rm kg/m^3}$ and $\delta_b \sim 1$:
$k_\mu \sim 0.1\,h\,{\rm Mpc}^{-1}$.  For linear perturbations
$\delta_b \ll 1$ at higher redshift, $k_\mu$ shifts to much larger $k$
(the entire observable range is in deep-MOND).

\paragraph{Grade: B (Derivation).}
The qualitative prediction (sign and mechanism) is derived rigorously
from the AQUAL equation with $\mu(x) = x/(1+x)$.  The quantitative
crossover scale requires mochi\_class.


\subsection{P26: Modified Power Spectrum Shape (Blue Tilt)}

\begin{prediction}[Power Spectrum Blue Tilt]
\label{pred:Pk}
The ratio $P^{\rm DFD}(k)/P^{\LCDM}(k)$ is a monotonically
decreasing function of $k$ in the range
$0.001 < k < 1\,h\,{\rm Mpc}^{-1}$, producing a ``blue tilt''
relative to \LCDM{} at large scales.
\end{prediction}

\paragraph{Mechanism.}  Deep-MOND growth enhances large-scale modes
more than small-scale modes (Finding~2 above).  The post-recombination
growth rate depends on $\Geff(k)$, which is larger for smaller $k$.

\paragraph{Implication for matter-radiation equality feature.}
The standard $P(k)$ turnover at $k_{\rm eq}$ arises because sub-horizon
modes during radiation domination undergo logarithmic growth (Meszaros
effect) while super-horizon modes grow normally.  In DFD, the baryonic
perturbations are tightly coupled to photons before decoupling (as in
\LCDM), so the transfer function at $z \sim 1100$ is similar.  But
the post-decoupling growth is scale-dependent, distorting the turnover
shape.  \textbf{The equality feature is not absent but modified:} its
amplitude and shape are altered by the $\Geff(k)$-dependent growth.

\paragraph{Grade: C (Conjecture).}  Qualitative direction robust;
quantitative amplitude requires mochi\_class.


\subsection{P27: Modified BAO Amplitude}

\begin{prediction}[BAO Broadband Amplitude Modification]
\label{pred:BAO}
The baryon acoustic oscillation (BAO) \emph{peak positions} are
unchanged (BAO is a geometric standard ruler set by the sound horizon
$r_d$, which DFD modifies by only $-0.10\%$, i15).  However, the
BAO \emph{amplitude} (contrast of the wiggles relative to the
broadband $P(k)$) is modified by the scale-dependent growth.  Deep-MOND
enhancement preferentially boosts the broadband relative to the wiggles,
reducing the BAO contrast by an estimated $\sim$10--30\% at $z < 1$.
\end{prediction}

\paragraph{Mechanism.}
BAO wiggles are imprinted at decoupling by acoustic oscillations in the
baryon-photon fluid.  After decoupling, in \LCDM, CDM provides
scale-independent potential wells and the wiggles are preserved with
amplitude set by $\Omega_b/\Omega_m$.  In DFD, the post-decoupling
growth is scale-dependent: modes at the BAO scale
($k_{\rm BAO} \sim 0.06\,h\,{\rm Mpc}^{-1}$) grow at a different rate
from adjacent modes.  The net effect is a \emph{partial washing-out}
of BAO wiggles relative to \LCDM, because the nonlinear $\mu$-dependent
growth does not respect the discrete $k$-structure of the acoustic
oscillations.

\paragraph{Grade: C (Conjecture).}  The sign is plausible but the
amplitude is uncertain.  BAO damping from nonlinear structure formation
occurs in \LCDM{} too (Eisenstein et al.\ 2007), so the DFD effect must
be disentangled from standard nonlinear BAO damping.  Requires
mochi\_class.


\subsection{P28: Early Galaxy Formation ($z > 6$) Enhanced}

\begin{prediction}[Enhanced Early Galaxy Formation]
\label{pred:early-gal}
DFD's deep-MOND growth rate $\delta \propto a^2$ predicts
earlier onset of nonlinear structure formation than \LCDM.
Structures go nonlinear ($\delta_b \sim 1$) at $z \sim 5$--10
in DFD, compared to $z \sim 1$--3 in \LCDM.  This naturally
accounts for the ``impossibly early galaxy'' problem from
JWST observations (Labbe et al.\ 2023, Castellano et al.\ 2022),
where massive galaxies appear at $z > 10$ in apparent tension
with \LCDM{} predictions.
\end{prediction}

\paragraph{Estimate.}  In deep-MOND, $\delta_b \propto a^2$, so
$\delta_b(z) = \delta_b(0)/(1+z)^2$.  For $\delta_b(0) \sim 1$
(present-day nonlinearity), $\delta_b \sim 1$ at $z = 0$.  But the
growth is faster: starting from the CMB normalization
$\delta_b(z=1100) \sim 10^{-5}$, growing as $a^2$ gives
$\delta_b \sim 10^{-5} \times (1100)^2 \sim 1$ at $z \sim 0$.
This matches the \LCDM{} result by construction (normalization).
However, the intermediate growth history differs: at $z = 10$,
$\delta \sim 10^{-5} \times (1100/11)^2 \sim 0.1$ in DFD vs
$\delta \sim 10^{-5} \times (1100/11) \sim 0.01$ in CDM.  So DFD
has $\delta$ an order of magnitude larger at $z = 10$ than \LCDM.
This allows halo formation to begin much earlier.

\paragraph{Caution.}  This is an order-of-magnitude argument.
The actual halo mass function requires the full nonlinear collapse
calculation in DFD, which has never been done.  Also, ``baryons only''
means no DM halo to form first --- baryonic collapse must be direct,
which changes the Jeans mass and cooling physics.

\paragraph{Grade: C (Conjecture).}  Qualitative direction correct;
quantitative halo mass function requires N-body or semi-analytic
calculation in the DFD framework.


%=============================================================================
\section{Stiff-Field Signatures Under $W'(0)=0$: Survival Assessment}
\label{sec:stiff}
%=============================================================================

The i19 stiff-field signatures were derived assuming $c_s = c$ for the
$\psi$-field.  The i20 $W'(0) = 0$ result modifies the picture:

\begin{itemize}
\item The \emph{temporal} kinetic term $\dot\psi^2/c^2$ is always
  present in the DFD action (it comes from the time-time component
  of the metric coupling, not from $W$).

\item The \emph{spatial} kinetic term from $W'(0)|\nabla\delta\psi|^2$
  vanishes at linear order around the deep-MOND background.

\item Therefore, for \textbf{free waves} (not sourced by matter), the
  equation $\ddot{\delta\psi} = 0$ (no spatial gradient restoring force)
  means \emph{the $\psi$-field does not support propagating waves in the
  deep-MOND regime}.  This is the ``$c_s^2 \to 0$ theorem'' from i15,
  now confirmed by the $W'(0)=0$ derivation.

\item For \textbf{sourced perturbations} (driven by baryonic matter
  through the AQUAL coupling), the $\psi$-field responds quasi-statically
  to the matter distribution: $\nabla\cdot[\mu(|\nabla\psi|/\astar)\nabla\psi]
  = -\text{source}$.  The field ``follows'' the matter rather than
  propagating independently.
\end{itemize}

\paragraph{Consequence for GW background (i19 signature (a)):}
The stiff-era GW enhancement depends on the background equation of state
of the $\psi$-field, not on perturbation dynamics.  If $\psi$ dominates
the energy density at some early epoch with $w = +1$ (stiff), primordial
GWs are amplified.  The question is whether the $\psi$-field actually
has $w = +1$ at early times.

With $W'(0) = 0$ and $W(y) \sim y^{3/2}$ for small $y$ (i20), the
kinetic energy density of a \emph{homogeneous} $\psi$-field is
$\rho_\psi \propto \dot\psi^2/c^2$ (temporal kinetic energy), while
the ``potential'' from $W$ depends on spatial gradients.  For a
homogeneous field, the spatial gradient vanishes and
$w = p/\rho = (\dot\psi^2/c^2)/(\dot\psi^2/c^2) = 1$ --- stiff.
\textbf{The stiff-era signature survives} if $\psi$ is ever the
dominant energy component.  However, the dust-branch analysis (i17)
shows $w \to 0$ (dust-like) for the cosmological $\psi$, suggesting
the stiff phase, if it exists, is brief and at very early times.

\textbf{Verdict: The blue-tilted GW background is POSSIBLE but
CONTINGENT on initial conditions, not a robust prediction.
Downgrade from B to C.}


%=============================================================================
\section{The Sharpest Near-Term Prediction}
\label{sec:sharpest}
%=============================================================================

\begin{tcolorbox}[colback=green!5!white, colframe=green!60!black,
  title=\textbf{SHARPEST NEAR-TERM DFD PREDICTION (2026--2027)}]

\textbf{P25: Tomographic $f\sigma_8(k)$ scale dependence.}

\textbf{Observable:} In redshift-space galaxy clustering, the
growth-rate parameter $f\sigma_8$ measured at different effective
wavenumber ranges.  DFD predicts $f\sigma_8$ is \emph{larger} at
$k \sim 0.02\,h\,{\rm Mpc}^{-1}$ than at $k \sim 0.2\,h\,{\rm Mpc}^{-1}$,
at fixed redshift.

\textbf{Why this is sharp:}
\begin{enumerate}
\item The \emph{sign} of the effect is derived from the AQUAL equation
  with zero free parameters: $\Geff \propto 1/\sqrt{g_N}$ means larger
  enhancement at larger scales.
\item The prediction is \emph{qualitative}: any detection of scale-dependent
  $f\sigma_8$ increasing toward large scales is a signature.  It does not
  require knowing the amplitude precisely.
\item \LCDM{} predicts strictly scale-independent $f\sigma_8$ at linear
  order.  Scale dependence in \LCDM{} arises only from nonlinear effects
  and is confined to $k > 0.1\,h\,{\rm Mpc}^{-1}$.
\item Superfluid DM and AeST both predict CDM-like scale-independent
  growth, making this a DFD-specific signature.
\end{enumerate}

\textbf{Test:}
\begin{itemize}
\item \textbf{Euclid DR1 (October 2026):} Spectroscopic galaxy sample
  in 4+ redshift bins.  Measure $f\sigma_8$ in broad $k$-bins
  ($k < 0.05$ vs $0.05 < k < 0.15$ vs $k > 0.15$).  Expected precision
  $\sim$5--10\% per bin.  A $\sim$30--50\% difference between low-$k$
  and high-$k$ $f\sigma_8$ (as DFD predicts in the deep-MOND regime)
  would be detectable at $>3\sigma$.
\item \textbf{DESI DR2 (2027):} Higher precision, more redshift bins.
  Can detect $\sim$10--20\% scale dependence at $>3\sigma$.
\end{itemize}

\textbf{Falsification criterion:} If Euclid DR1 + DESI DR2 both
find $f\sigma_8$ scale-independent to $<10\%$ precision across
$0.02 < k < 0.3\,h\,{\rm Mpc}^{-1}$ at all redshifts $0.5 < z < 2$,
the deep-MOND growth regime is constrained to apply only below the
survey's minimum $k$, which would severely limit DFD's ability to
produce enough structure.
\end{tcolorbox}


%=============================================================================
\section{Revised Prediction Audit}
\label{sec:audit}
%=============================================================================

\begin{longtable}{p{0.5cm}p{4cm}p{1cm}p{1cm}p{5.5cm}}
\caption{Full prediction audit, updated to i21.} \\
\toprule
\textbf{\#} & \textbf{Prediction} & \textbf{i19} & \textbf{i21}
  & \textbf{Notes} \\
\midrule
\endfirsthead
\toprule
\textbf{\#} & \textbf{Prediction} & \textbf{i19} & \textbf{i21}
  & \textbf{Notes} \\
\midrule
\endhead
T0 & GWs never grav.\ lensed & A & A & Structural theorem. Unchanged. \\
P2 & Shadow $+4.6\%$ & A & A & Algebraic. Unchanged. \\
P5 & $c_T = c$ exactly & A & A & Structural. Unchanged. \\
P20 & SNe Ia $\psi$-screen anisotropy & A & A & Distance-bias. Unchanged. \\
P21 & $D_L^{\rm EM}/D_L^{\rm GW}$ ratio & A & A & Structural. Unchanged. \\
P22 & FRB DM-$z$ excess $+7$--12\% & A & A & Distance-bias. Unchanged. \\
P24 & $\beta = 0$ exactly & A & A & Microsector S$^3$ CS. Unchanged. \\
& & & & True tension $\sim$2--3$\sigma$ (i20). \\
P12 & Zero scalar GW memory & A & A & PTA breathing mode. Unchanged. \\
P18 & IFMR radial gradient & A & A & Unchanged. \\
\midrule
P1 & Lunar nonreciprocal 0.93~ns & B & B & Unchanged. \\
P3 & BAO positions (same $r_d$) & B & B & $r_d$ modified $-0.10\%$. Unchanged. \\
P4 & $\Delta H_0 = 3.9 \pm 1.1$ km/s/Mpc & B & B & Distance-bias. Unchanged. \\
P7 & Q-ratio 0.52$\pm$0.08 & B & B & Unchanged. \\
P8 & Void lensing 1.5--3.5$\times$ & B & B & Unchanged. \\
P9 & Tomographic $S_8$ gradient & B & B & Survives $c_s=c$ (i19). \\
P10 & A$_L$ scale dependence & B & B & Unchanged. \\
P11 & kSZ velocity $+10\%$ & B & B & Unchanged. \\
P19 & TDE rate $+40\%$ in LSB gal.\ & B & B & Unchanged. \\
P25 & $f\sigma_8(k,z)$ scale dep.\ & --- & B & \textbf{NEW}. Finding 1. \\
\midrule
P6 & $\sigma_8$ value & B & C & DOWNGRADED. Range 0.5--2. \\
P13 & $\Sigma m_\nu = 61.4$ meV & C & C & Unchanged. 2.0 meV margin. \\
P14 & Excess large voids $+24\%$ & C & C & Unchanged. \\
P15 & $\alpha = 1/137$ & C & C & Unchanged (microsector). \\
P16 & Bar pattern speed $+19\%$ & C & C & Overclaimed (MOND, not DFD-specific). \\
P17 & 21cm bispectrum & C & C & Unchanged. \\
P23 & Stiff-era GW blue tilt & B & C & DOWNGRADED. Contingent on \\
& & & & initial conditions (Finding 3). \\
P26 & $P(k)$ blue tilt & --- & C & \textbf{NEW}. Finding 2. \\
P27 & BAO amplitude reduction & --- & C & \textbf{NEW}. Sec.~2.3. \\
P28 & Enhanced early galaxies & --- & C & \textbf{NEW}. Sec.~2.4. \\
\bottomrule
\end{longtable}

\textbf{Updated totals:} 9 Theorem-grade (A), 10 Derivation-grade (B),
10 Conjecture (C).  29 predictions total.  Honest (A+B):parameter
ratio = 19:1.

\textbf{Key changes from i19:}
\begin{itemize}
\item $\sigma_8$ value downgraded B $\to$ C (factor $\sim$4 range).
\item Stiff-era GW blue tilt downgraded B $\to$ C (contingent on ICs).
\item 4 new predictions added (P25--P28), all from deep-MOND cosmology.
\item P25 (scale-dependent $f\sigma_8$) is the most observationally
  accessible new prediction.
\end{itemize}


%=============================================================================
\section*{Derivation Chain Summary}
%=============================================================================

Every claim traces to:
\begin{enumerate}
\item \textbf{i20 Finding 1} (Degenerate Linearization Theorem): $W'(0) = 0$
  for $\mu(x) = x/(1+x)$. Linearized PT fails.
\item \textbf{i20 Finding 2} ($W'(0) = 0$ is correct, not $W'(0) = 1$):
  Explicit calculation of $W(y)$ for simple $\mu$.
\item \textbf{i20 Finding 3} (Nonlinear corrections, deep-MOND):
  AQUAL perturbation theory following Nusser (2002).
\item \textbf{i20 Finding 4} ($\Geff$ two-regime structure):
  $G_N$ at linear order, $G_N\sqrt{\aM/g_N}$ at nonlinear AQUAL order.
\item \textbf{i19 Prediction Audit}: 9A, 10B, 6C baseline.
\item \textbf{i19 Stiff-field signatures}: $c_s = c$, three signatures.
\item \textbf{i20 $\beta = 0$ assessment}: True tension $\sim$2--3$\sigma$.
\item \textbf{i17 Dust-branch analysis}: $w \to 0$ for cosmological $\psi$.
\item \textbf{Hwang \& Noh (2025)}: Independent confirmation of deep-MOND
  growth $\delta \propto a^2$ (arXiv:2410.10205).
\item \textbf{DFD v3.2}: Action (Eq.~2.14), microsector (Appendices L, N),
  $\mu$-function (Section 2).
\end{enumerate}

\end{document}
